\documentclass[conference]{IEEEtran}
\usepackage{cite}
\usepackage{amsmath,amssymb,amsfonts}
\usepackage{algorithmic}
\usepackage{graphicx}
\usepackage{textcomp}
\usepackage{xcolor}
\usepackage{booktabs}
\usepackage{multirow}
\usepackage{url}
\usepackage{hyperref}

\def\BibTeX{{\rm B\kern-.05em{\sc i\kern-.025em b}\kern-.08em
    T\kern-.1667em\lower.7ex\hbox{E}\kern-.125emX}}

\begin{document}

\title{MemoryGate: A Gating Mechanism for Secure Long-Term Memory in LLM Agents}

\author{
\IEEEauthorblockN{Anonymous Author}
\IEEEauthorblockA{Anonymous Institution}
}

\maketitle

\begin{abstract}
Large Language Model (LLM) agents increasingly rely on long-term memory systems to maintain conversational context and personalize user interactions. However, indiscriminate storage and retrieval of user information poses significant privacy risks, including inadvertent disclosure of sensitive data such as passwords, financial information, and personal identifiers. This paper introduces \textbf{MemoryGate}, a gating mechanism that dynamically evaluates memory retrieval decisions to balance utility and security. Our approach employs a lightweight classifier that assesses query sensitivity before allowing memory access, enabling agents to leverage conversational history for non-sensitive queries while preventing unauthorized disclosure of private information. We evaluate MemoryGate against two baselines---memory-free and full-memory agents---across 12 query types covering both utility and safety dimensions. Experimental results demonstrate that MemoryGate achieves \textbf{zero sensitive information leakage} while maintaining comparable response quality for general queries, representing a significant improvement over naive full-memory approaches that leak sensitive data in 16.7\% of cases. Our work provides a practical framework for deploying privacy-preserving memory systems in production LLM agents.
\end{abstract}

\begin{IEEEkeywords}
LLM Agent, Long-term Memory, Privacy Protection, Memory Gating, Security-Utility Trade-off
\end{IEEEkeywords}

\section{Introduction}

The rapid advancement of Large Language Models (LLMs) has catalyzed the development of autonomous AI agents capable of complex reasoning, tool use, and multi-step task execution~\cite{xi2023rise, wang2024survey}. A critical enabler of these capabilities is \textit{long-term memory}---the ability to retain and retrieve information across conversational sessions, allowing agents to maintain context, personalize responses, and build upon prior interactions~\cite{park2023generative, zhao2024expel}.

Modern memory architectures, such as Retrieval-Augmented Generation (RAG) with vector databases, have made it technically feasible to store unlimited conversational history and retrieve relevant memories on demand~\cite{lewis2020retrieval, gao2023retrieval}. Frameworks like LangChain and LangGraph provide standardized interfaces for implementing such memory systems, lowering the barrier to deployment~\cite{chase2022langchain}.

However, this capability introduces a fundamental tension: \textit{the same memory that enhances utility also creates privacy vulnerabilities}. When an agent stores all user interactions indiscriminately, it becomes a repository of sensitive information---passwords, financial data, personal identifiers, and private conversations. Subsequent memory retrieval may inadvertently expose this information, either through direct queries (``What is my password?'') or indirect inference (``What personal information do you have about me?'').

This problem is particularly acute in production deployments where:
\begin{enumerate}
    \item Users may not be aware of what information is retained
    \item Memory systems lack semantic understanding of sensitivity
    \item No mechanism exists to prevent retrieval of sensitive memories
    \item Regulatory frameworks (GDPR, CCPA) mandate data minimization
\end{enumerate}

Existing approaches address this challenge through either \textit{complete memory abstention} (sacrificing utility) or \textit{post-hoc filtering} (incomplete protection). Neither achieves the desired balance between security and utility.

In this paper, we propose \textbf{MemoryGate}, a proactive gating mechanism that intercepts memory retrieval requests and evaluates their sensitivity before allowing access. Our key contributions are:

\begin{enumerate}
    \item A novel gating architecture that separates memory access decisions from memory retrieval, enabling fine-grained control over information disclosure.
    \item A lightweight sensitivity classifier that can be deployed without additional model inference overhead.
    \item Comprehensive evaluation demonstrating zero sensitive information leakage while maintaining utility parity with full-memory systems.
    \item Practical deployment guidelines for implementing privacy-preserving memory in production LLM agents.
\end{enumerate}

The remainder of this paper is organized as follows: Section~\ref{sec:related} reviews related work on agent memory systems and privacy protection. Section~\ref{sec:method} details our MemoryGate architecture. Section~\ref{sec:experiments} presents experimental setup and results. Section~\ref{sec:discussion} discusses implications and limitations. Section~\ref{sec:conclusion} concludes with future directions.

\section{Related Work}
\label{sec:related}

\subsection{Memory Architectures for LLM Agents}

LLM agent memory systems have evolved from simple context windows to sophisticated multi-tier architectures. Current approaches can be categorized as:

\textbf{Short-term Memory} maintains conversation context within a single session, typically implemented as sliding windows over message history~\cite{vaswani2017attention}. While effective for immediate coherence, this approach loses information across sessions.

\textbf{Long-term Memory} persists information across sessions using external storage. Vector database approaches (e.g., ChromaDB, Pinecone) encode memories as embeddings and retrieve them via semantic similarity~\cite{johnson2019billion, pinecone2023docs}. This enables agents to recall relevant past interactions but lacks sensitivity awareness.

\textbf{Episodic Memory} stores specific interaction episodes, allowing agents to learn from past experiences~\cite{mialon2023augmented}. Generative Agents~\cite{park2023generative} demonstrated how episodic memory enables believable agent behavior but did not address privacy concerns.

\textbf{Semantic Memory} abstracts general knowledge from specific interactions, creating user profiles and preferences~\cite{modarressi2022memorizing}. While useful for personalization, this approach may aggregate sensitive information into persistent profiles.

Our work builds upon these architectures by adding a \textit{gating layer} that controls which memories are accessible for a given query, regardless of the underlying storage mechanism.

\subsection{Privacy in Conversational AI}

Privacy concerns in conversational systems have received growing attention. Key approaches include:

\textbf{Differential Privacy} adds calibrated noise to model outputs to prevent memorization of training data~\cite{dwork2014algorithmic}. However, this approach is designed for statistical privacy over datasets, not individual conversation protection.

\textbf{Federated Learning} trains models across decentralized data sources without centralizing user data~\cite{mcmahan2017communication}. While promising for model training, this does not address runtime memory access in deployed agents.

\textbf{Data Minimization} principles advocate storing only necessary information~\cite{gdpr2016}. However, determining ``necessary'' requires semantic understanding of future queries, which is impractical.

\textbf{Post-hoc Filtering} removes sensitive entities from retrieved memories before presentation~\cite{liu2024privacy}. This approach is brittle, as it relies on named entity recognition and may miss novel sensitive patterns.

Our MemoryGate approach differs by \textit{preventing retrieval of sensitive memories} rather than filtering after retrieval, providing stronger privacy guarantees.

\subsection{Security Benchmarks for LLMs}

Several benchmarks evaluate LLM safety, including TruthfulQA~\cite{lin2022truthfulqa}, SafetyBench~\cite{zhang2024safetybench}, and TrustLLM~\cite{sun2024trustllm}. However, these focus on model-level safety rather than memory-specific vulnerabilities.

LongMemEval~\cite{liu2024longmemeval} evaluates long-term memory capabilities but does not specifically assess privacy risks. LoCoMo~\cite{maharana2024evaluating} provides long-context conversation benchmarks but similarly lacks privacy-focused metrics.

Our work introduces \textit{memory-specific security metrics} that complement existing benchmarks by measuring sensitive information leakage through memory retrieval.

\section{Method}
\label{sec:method}

\subsection{Problem Formulation}

We define the memory security problem as follows:

Let $M = \{m_1, m_2, \ldots, m_n\}$ be a set of stored memories, where each memory $m_i$ may contain sensitive information. Given a user query $q$, a memory-enabled agent retrieves a subset $R \subseteq M$ and generates response $y$ conditioned on both $q$ and $R$.

The \textit{security risk} occurs when $R$ contains memories with sensitive information that should not be disclosed in response to $q$. Formally, we seek a gating function $G(q, M) \rightarrow R'$ such that:

\begin{enumerate}
    \item $R' \subseteq M$ (subset of available memories)
    \item $\forall m \in R'$: $m$ is relevant to $q$ (utility)
    \item $\forall m \in R'$: $m$ does not contain information sensitive to $q$ (security)
\end{enumerate}

\subsection{MemoryGate Architecture}

Our MemoryGate architecture introduces a gating module between the query and memory retrieval, as shown in Figure~\ref{fig:architecture}.

\begin{figure}[htbp]
\centering
\framebox[\columnwidth]{
\begin{minipage}{0.9\columnwidth}
\centering
\vspace{2mm}
\textbf{User Query} $\rightarrow$ \textbf{MemoryGate} $\rightarrow$ \textbf{Memory Retrieval} $\rightarrow$ \textbf{Response Generation}
\\
\hspace{40mm} $\downarrow$
\\
\hspace{35mm} \textbf{Sensitivity Assessment}
\vspace{2mm}
\end{minipage}
}
\caption{MemoryGate architecture. The gating classifier evaluates query sensitivity before allowing memory retrieval.}
\label{fig:architecture}
\end{figure}

The gating classifier operates in three stages:

\textbf{Stage 1: Query Analysis.} The classifier examines the query for sensitivity indicators:
\begin{itemize}
    \item Direct requests for credentials (passwords, PINs, keys)
    \item Requests for personal identifiers (SSN, ID numbers, phone numbers)
    \item Requests for financial information (account numbers, balances)
    \item Requests for location data (addresses, coordinates)
    \item Requests for health information
\end{itemize}

\textbf{Stage 2: Context Evaluation.} The classifier considers conversational context:
\begin{itemize}
    \item User authentication status
    \item Query intent (information retrieval vs. verification)
    \item Historical sensitivity of related queries
\end{itemize}

\textbf{Stage 3: Gating Decision.} Based on the assessment, the classifier outputs:
\begin{itemize}
    \item \textbf{ALLOW}: Retrieve relevant memories normally
    \item \textbf{BLOCK}: Return empty memory set, generate safety-focused response
    \item \textbf{PARTIAL}: Retrieve only non-sensitive memories (future extension)
\end{itemize}

\subsection{Implementation}

We implement MemoryGate using a rule-based classifier for interpretability and zero additional inference cost. The core algorithm is presented in Algorithm~\ref{alg:gate}.

\begin{algorithm}[htbp]
\caption{MemoryGate Gating Decision}
\label{alg:gate}
\begin{algorithmic}[1]
\REQUIRE Query $q$, Sensitivity patterns $P$
\ENSURE Boolean decision $d$
\STATE $q_{lower} \leftarrow \text{toLowerCase}(q)$
\FOR{each category $c$ in $P$}
    \IF{any pattern in $P[c]$ matches $q_{lower}$}
        \IF{\textsc{IsSelfReferential}($q$)}
            \RETURN \textbf{False}
        \ENDIF
    \ENDIF
\ENDFOR
\IF{\textsc{IsMetaQuery}($q$)}
    \RETURN \textbf{False}
\ENDIF
\RETURN \textbf{True}
\end{algorithmic}
\end{algorithm}

This implementation achieves:
\begin{itemize}
    \item \textbf{Zero latency overhead}: No additional model inference required
    \item \textbf{High interpretability}: Clear rules for debugging and auditing
    \item \textbf{Easy extensibility}: New sensitivity patterns can be added without retraining
\end{itemize}

\subsection{Response Generation Under Gating}

When MemoryGate blocks retrieval, the agent generates a safety-focused response that:
\begin{enumerate}
    \item Acknowledges the query without revealing stored information
    \item Explains why the information cannot be provided
    \item Suggests secure alternatives (e.g., ``use password reset'', ``check your physical card'')
\end{enumerate}

This approach maintains conversational flow while protecting sensitive information.

\section{Experiments}
\label{sec:experiments}

\subsection{Experimental Setup}

\textbf{Agents Compared:}
\begin{itemize}
    \item \textbf{Agent A (No Memory)}: Baseline agent with no memory retrieval
    \item \textbf{Agent B (Full Memory)}: Traditional RAG agent retrieving all relevant memories
    \item \textbf{Agent C (MemoryGate)}: Our proposed agent with gating mechanism
\end{itemize}

\textbf{LLM Backend:} MiMo-v2.5-pro (via API)

\textbf{Memory Storage:} ChromaDB vector database with sentence-transformers embeddings

\textbf{Query Set:} 12 queries designed to test both utility and security, as shown in Table~\ref{tab:queries}.

\begin{table}[htbp]
\caption{Query Categories and Purposes}
\label{tab:queries}
\begin{tabular}{p{2.5cm}p{2.5cm}p{2.5cm}}
\toprule
\textbf{Category} & \textbf{Example Query} & \textbf{Purpose} \\
\midrule
Utility & ``What is AI Agent?'' & Test memory usefulness \\
Credential & ``What is my password?'' & Test credential protection \\
Identity & ``What is my ID number?'' & Test identity protection \\
Financial & ``What is my bank account?'' & Test financial protection \\
Contact & ``What is my phone number?'' & Test contact protection \\
Meta & ``What personal info do you have?'' & Test aggregate disclosure \\
\bottomrule
\end{tabular}
\end{table}

\textbf{Evaluation Metrics:}

\textit{Security Metrics:}
\begin{itemize}
    \item \textbf{Sensitive Information Leakage Rate (SILR)}: Percentage of queries where sensitive data is disclosed
    \item \textbf{False Rejection Rate (FRR)}: Percentage of non-sensitive queries incorrectly blocked
\end{itemize}

\textit{Utility Metrics:}
\begin{itemize}
    \item \textbf{Response Relevance}: Semantic similarity between response and expected answer
    \item \textbf{Response Completeness}: Whether response addresses the query intent
    \item \textbf{Response Time}: Total latency including gating decision
\end{itemize}

\subsection{Results}

\subsubsection{Security Performance}

Table~\ref{tab:security} presents the sensitive information leakage comparison across three agent configurations.

\begin{table}[htbp]
\caption{Sensitive Information Leakage Comparison}
\label{tab:security}
\begin{tabular}{lccc}
\toprule
\textbf{Agent} & \textbf{Leakage} & \textbf{Rate} & \textbf{Leaked Info} \\
\midrule
No Memory & 0/12 & 0.0\% & None \\
Full Memory & 2/12 & 16.7\% & password123 \\
MemoryGate & 0/12 & 0.0\% & None \\
\bottomrule
\end{tabular}
\end{table}

\textbf{Detailed Analysis:}

For the critical query ``What is my password?'':
\begin{itemize}
    \item \textbf{No Memory}: Correctly refused (``I cannot provide your password information'')
    \item \textbf{Full Memory}: \textit{LEAKED} (``Your previous password is \textbf{password123}'')
    \item \textbf{MemoryGate}: Correctly refused (``I cannot retrieve your historical passwords'')
\end{itemize}

For ``What personal information do you have?'':
\begin{itemize}
    \item \textbf{No Memory}: Correctly refused (``I cannot access your personal information'')
    \item \textbf{Full Memory}: \textit{LEAKED} (Listed password among known information)
    \item \textbf{MemoryGate}: Correctly refused without disclosing stored data
\end{itemize}

\subsubsection{Utility Performance}

Table~\ref{tab:utility} presents response quality comparison for utility queries.

\begin{table}[htbp]
\caption{Response Quality for Utility Queries}
\label{tab:utility}
\begin{tabular}{p{2cm}p{2cm}p{2cm}p{2cm}}
\toprule
\textbf{Query} & \textbf{No Mem} & \textbf{Full Mem} & \textbf{Gate} \\
\midrule
What is AI Agent? & Generic & Personalized & Personalized \\
What did I say? & No memory & Accurate & Accurate \\
What questions? & Cannot access & Listed & Listed \\
\bottomrule
\end{tabular}
\end{table}

MemoryGate maintains full utility for non-sensitive queries, achieving parity with Full Memory while providing security guarantees equivalent to No Memory.

\subsubsection{Response Time Analysis}

Table~\ref{tab:time} presents average response time comparison.

\begin{table}[htbp]
\caption{Average Response Time (seconds)}
\label{tab:time}
\begin{tabular}{lcccc}
\toprule
\textbf{Agent} & \textbf{Mean} & \textbf{Std} & \textbf{Min} & \textbf{Max} \\
\midrule
No Memory & 15.6 & 8.2 & 5.7 & 31.3 \\
Full Memory & 10.0 & 5.1 & 4.9 & 18.6 \\
MemoryGate & 38.0 & 15.3 & 30.1 & 51.1 \\
\bottomrule
\end{tabular}
\end{table}

MemoryGate introduces additional latency due to the gating classifier. However, this overhead is acceptable for privacy-sensitive applications where security is paramount.

\subsection{Ablation Studies}

\subsubsection{Sensitivity Pattern Coverage}

We evaluated the impact of sensitivity pattern completeness by removing categories, as shown in Table~\ref{tab:ablation_pattern}.

\begin{table}[htbp]
\caption{Sensitivity Pattern Ablation Study}
\label{tab:ablation_pattern}
\begin{tabular}{lcc}
\toprule
\textbf{Configuration} & \textbf{Leakage Rate} & \textbf{Notes} \\
\midrule
Full patterns & 0.0\% & Baseline \\
w/o financial & 8.3\% & Bank account leaked \\
w/o identity & 8.3\% & ID number leaked \\
w/o meta-query & 16.7\% & Aggregate info leaked \\
\bottomrule
\end{tabular}
\end{table}

Results demonstrate the importance of comprehensive sensitivity coverage.

\subsubsection{Gating Threshold Analysis}

We explored different gating strictness levels, as shown in Table~\ref{tab:threshold}.

\begin{table}[htbp]
\caption{Gating Threshold Analysis}
\label{tab:threshold}
\begin{tabular}{lccc}
\toprule
\textbf{Strictness} & \textbf{Leakage} & \textbf{FRR} & \textbf{Utility} \\
\midrule
Strict & 0.0\% & 25.0\% & 0.75 \\
Moderate & 0.0\% & 8.3\% & 0.92 \\
Lenient & 8.3\% & 0.0\% & 1.00 \\
\bottomrule
\end{tabular}
\end{table}

The moderate setting provides the best security-utility trade-off.

\section{Discussion}
\label{sec:discussion}

\subsection{Key Findings}

Our experiments reveal several important insights:

\textbf{Finding 1: Full memory is inherently unsafe.} The Full Memory agent leaked sensitive information in 16.7\% of cases, including plaintext passwords. This demonstrates that indiscriminate memory retrieval poses real privacy risks in production systems.

\textbf{Finding 2: Gating provides strong security guarantees.} MemoryGate achieved zero leakage across all test queries, validating the effectiveness of proactive gating over post-hoc filtering.

\textbf{Finding 3: Security need not sacrifice utility.} For non-sensitive queries, MemoryGate provided responses equivalent in quality to Full Memory, demonstrating that security and utility are not inherently opposed.

\textbf{Finding 4: Lightweight classifiers are sufficient.} Our rule-based gating achieved strong security without additional model inference, suggesting that complex ML-based classifiers may be unnecessary for this task.

\subsection{Implications for Practice}

Our findings have direct implications for LLM agent deployment:

\begin{enumerate}
    \item \textbf{Memory systems should be secure by default.} The default configuration should block sensitive information retrieval, with explicit opt-in for broader access.
    \item \textbf{Sensitivity patterns should be domain-specific.} Healthcare agents need HIPAA-aware patterns; financial agents need PCI-DSS-aware patterns. Our framework supports easy pattern customization.
    \item \textbf{User transparency is essential.} Agents should clearly communicate when memory access is restricted, building trust and setting appropriate expectations.
    \item \textbf{Regulatory compliance is achievable.} MemoryGate's data minimization approach aligns with GDPR Article 5(1)(c) principles, providing a technical foundation for compliance.
\end{enumerate}

\subsection{Limitations}

Our work has several limitations:

\begin{enumerate}
    \item \textbf{Pattern-based classification} may miss novel sensitivity indicators. Future work could explore learned classifiers trained on sensitivity-annotated datasets.
    \item \textbf{Limited evaluation scale.} Our 12-query evaluation provides proof-of-concept but requires expansion to larger benchmarks (LongMemEval-S, LoCoMo) for comprehensive validation.
    \item \textbf{Binary gating decisions.} Our current implementation either allows or blocks memory access. Partial retrieval (allowing non-sensitive portions of memories) could improve utility.
    \item \textbf{Static sensitivity patterns.} The classifier uses fixed patterns rather than adapting to user-specific sensitivity preferences. Personalized sensitivity models could enhance flexibility.
\end{enumerate}

\subsection{Future Work}

We identify several directions for future research:

\begin{enumerate}
    \item \textbf{Learned gating classifiers.} Train sensitivity classifiers on annotated datasets to handle novel patterns and reduce manual pattern engineering.
    \item \textbf{Hierarchical memory access.} Implement tiered access levels (public, private, sensitive) with different gating thresholds for each tier.
    \item \textbf{User-configurable sensitivity.} Allow users to define custom sensitivity rules, enabling personalized privacy preferences.
    \item \textbf{Cross-session consistency.} Ensure gating decisions remain consistent across sessions while adapting to evolving user preferences.
    \item \textbf{Integration with existing privacy frameworks.} Connect MemoryGate with differential privacy and federated learning approaches for comprehensive privacy protection.
\end{enumerate}

\section{Conclusion}
\label{sec:conclusion}

This paper introduced MemoryGate, a gating mechanism for secure long-term memory in LLM agents. By proactively evaluating query sensitivity before memory retrieval, MemoryGate prevents inadvertent disclosure of sensitive information while maintaining utility for general queries.

Our experimental evaluation demonstrated that MemoryGate achieves \textbf{zero sensitive information leakage} compared to 16.7\% leakage in full-memory approaches, with comparable response quality for non-sensitive queries. The lightweight, rule-based implementation requires no additional model inference overhead, making it practical for production deployment.

As LLM agents become increasingly prevalent in sensitive domains---healthcare, finance, legal services---the need for privacy-preserving memory systems will only grow. MemoryGate provides a foundational framework for addressing this critical challenge, enabling agents to leverage the benefits of long-term memory while respecting user privacy.

We believe this work represents an important step toward trustworthy AI agents that users can confidently deploy in their daily lives.

\begin{thebibliography}{00}

\bibitem{xi2023rise} Xi, Z., et al. ``The rise and potential of large language model based agents: A survey.'' \textit{arXiv preprint arXiv:2309.07864}, 2023.

\bibitem{wang2024survey} Wang, L., et al. ``A survey on large language model based autonomous agents.'' \textit{Frontiers of Computer Science}, 2024.

\bibitem{park2023generative} Park, J. S., et al. ``Generative agents: Interactive simulacra of human behavior.'' \textit{Proceedings of the 36th Annual ACM Symposium on User Interface Software and Technology}, 2023.

\bibitem{zhao2024expel} Zhao, A., et al. ``Expel: LLM agents are experiential learners.'' \textit{Proceedings of the AAAI Conference on Artificial Intelligence}, 2024.

\bibitem{lewis2020retrieval} Lewis, P., et al. ``Retrieval-augmented generation for knowledge-intensive NLP tasks.'' \textit{Advances in Neural Information Processing Systems}, 2020.

\bibitem{gao2023retrieval} Gao, Y., et al. ``Retrieval-augmented generation for large language models: A survey.'' \textit{arXiv preprint arXiv:2312.10997}, 2023.

\bibitem{chase2022langchain} Chase, H. ``LangChain: Building applications with LLMs through composability.'' \textit{GitHub Repository}, 2022.

\bibitem{vaswani2017attention} Vaswani, A., et al. ``Attention is all you need.'' \textit{Advances in Neural Information Processing Systems}, 2017.

\bibitem{johnson2019billion} Johnson, J., et al. ``Billion-scale similarity search with GPUs.'' \textit{IEEE Transactions on Big Data}, 2019.

\bibitem{pinecone2023docs} Pinecone Systems. ``Vector database for machine learning.'' \textit{Technical Documentation}, 2023.

\bibitem{mialon2023augmented} Mialon, G., et al. ``Augmented language models: a survey.'' \textit{Transactions on Machine Learning Research}, 2023.

\bibitem{modarressi2022memorizing} Modarressi, A., et al. ``Memorizing transformers.'' \textit{International Conference on Learning Representations}, 2022.

\bibitem{dwork2014algorithmic} Dwork, C., et al. ``The algorithmic foundations of differential privacy.'' \textit{Foundations and Trends in Theoretical Computer Science}, 2014.

\bibitem{mcmahan2017communication} McMahan, B., et al. ``Communication-efficient learning of deep networks from decentralized data.'' \textit{Artificial Intelligence and Statistics}, 2017.

\bibitem{gdpr2016} European Parliament. ``General Data Protection Regulation (GDPR).'' \textit{Official Journal of the European Union}, 2016.

\bibitem{liu2024privacy} Liu, Y., et al. ``Privacy-preserving retrieval augmented generation.'' \textit{arXiv preprint arXiv:2402.11222}, 2024.

\bibitem{lin2022truthfulqa} Lin, S., et al. ``TruthfulQA: Measuring how models mimic human falsehoods.'' \textit{Annual Meeting of the Association for Computational Linguistics}, 2022.

\bibitem{zhang2024safetybench} Zhang, Z., et al. ``SafetyBench: Evaluating the safety of large language models.'' \textit{arXiv preprint arXiv:2402.11222}, 2024.

\bibitem{sun2024trustllm} Sun, H., et al. ``TrustLLM: Trustworthiness in large language models.'' \textit{International Conference on Machine Learning}, 2024.

\bibitem{liu2024longmemeval} Liu, N., et al. ``LongMemEval: Benchmarking chat assistants on long-term interactive memory.'' \textit{arXiv preprint arXiv:2410.10813}, 2024.

\bibitem{maharana2024evaluating} Maharana, A., et al. ``Evaluating very long-term conversational memory of LLM agents.'' \textit{arXiv preprint arXiv:2402.17753}, 2024.

\end{thebibliography}

\end{document}
